% =============================================================================
% CHAPTER 1 — INTRODUCTION GÉNÉRALE
% =============================================================================

\chapter*{Introduction Générale}
\addcontentsline{toc}{chapter}{Introduction Générale}
\markboth{Introduction Générale}{}
\label{chap:introduction}

Dans un contexte marqué par la transformation numérique et l'émergence rapide des technologies, le marché de l'emploi connaît des mutations profondes qui redéfinissent les compétences attendues par les sociétés. Cette évolution génère un écart entre les compétences enseignées dans les établissements académiques et les exigences réelles des métiers du transport et de la logistique. Face à ces changements, les établissements d'enseignement supérieur sont confrontés à un défi majeur qui consiste à assurer une adéquation entre les compétences acquises par les étudiants et les exigences réelles du marché du travail. 

\noindent Malgré la qualité des formations dispensées, de nombreux étudiants rencontrent des difficultés lors de leur insertion professionnelle en raison d'un manque de visibilité sur leur niveau réel d'employabilité, les compétences recherchées par les recruteurs et les écarts existants entre leur profil et les métiers ciblés. Parallèlement, les responsables pédagogiques disposent rarement d'outils leur permettant d'évaluer objectivement l'adéquation entre les formations académiques et les besoins du marché de l'emploi. De leur côté, les sociétés font face à un volume croissant de candidatures, rendant complexe l'identification rapide des profils les plus pertinents. 

\noindent Ce projet de fin d’études, réalisé dans le cadre du Programme d'Encouragement des Jeunes Chercheurs 2025 (PEJC), en collaboration entre l'Institut Supérieur de Gestion Industrielle de Sfax (ISGIS) et la Fondation Djagora. L'objectif principal de ce projet est de concevoir et de développer LogiLink, une plateforme IA d'adéquation Formation–Emploi pour les métiers Transport \& Logistique.

\noindent Dans le cadre de ce projet réalisé en équipe, notre contribution s'est principalement concentrée sur la conception et le développement des modules de génération de CV, calcul du score ATS et employabilité. Plus précisément, nous avons participé au développement du moteur de génération automatisée de CV, utilisé comme base pour le mécanisme de calcul du score ATS permettant d'évaluer le profil par rapport à un poste visé. En parallèle, nous avons contribué à la conception du modèle de calcul du score d'employabilité fondé sur l'exploitation d'une matrice de couverture associant les matières aux différents métiers. Enfin, nous avons participé au développement du modèle d'employabilité prévisionnelle destiné à estimer prévoir l'évolution potentielle de l'employabilité d'un étudiant en fonction des  recommandations générées par la plateforme et validées par les responsables pédagogiques.

\vspace{12pt}

\noindent Afin de restituer de manière claire et structurée la démarche adoptée, ce rapport est organisé en trois chapitres distincts :

\begin{itemize}[leftmargin=*]
  \item \textbf{Chapitre 1 : Présentation et Analyse du Système}
  
  \noindent Le premier chapitre est consacré à la présentation du cadre général du projet et de l'organisme d'accueil, afin de poser le contexte de notre étude. Il expose ensuite la problématique, qui sert de point de départ à une analyse critique de l'existant pour introduire la solution. Enfin, ce chapitre détaille la spécification des exigences fonctionnelles et non fonctionnelles, tout en formalisant les besoins à travers le Product Backlog établi selon la méthodologie Scrum.
  \item \textbf{Chapitre 2 : Conception du système} 

  \noindent Le deuxième chapitre est consacré à la conception du système. Il décrit l'architecture de la plateforme ainsi que la modélisation UML adoptée pour formaliser les différents aspects fonctionnels et techniques de la solution.
  \item \textbf{Chapitre 3 : Realisation du systeme} 
  
  \noindent Le troisième chapitre porte sur la réalisation de la plateforme. Il présente l'environnement de développement, les technologies adoptées ainsi que les principaux modules implémentés, notamment ceux relatifs à la génération de CV, au calcul du score ATS, du score d'employabilité et d'employabilité prévisionnelle. Enfin, il met en évidence les interfaces développées et les résultats obtenus.
\end{itemize}

\vspace{12pt}
